\plannedchapter{构建流程、编译参数和工具链心智模型}
{本章把 C/C++ 的实际构建流程讲清楚：源文件如何被组织，编译命令如何影响语义和性能，CMake、编译数据库和构建类型在工程中承担什么角色。}
{\item 为什么同一份源码在不同 flags 下会生成完全不同的机器码？
\item Debug、Release、RelWithDebInfo、sanitizer、coverage build 分别解决什么问题？
\item 什么是 translation unit、include path、library path、compile command database？}
{\item 编译驱动程序如何调用预处理器、编译器、汇编器和链接器。
\item 优化级别、目标 ISA、调试信息、异常/RTTI、LTO 对二进制的影响。
\item 构建系统如何决定哪些文件重新编译，为什么大型工程需要明确依赖。}
{建立一个多文件 C++ 小工程，用 GCC 和 Clang 分别构建；生成 \code{compile_commands.json}；比较 \code{-O0}、\code{-O3}、\code{-march=native}、\code{-g} 的汇编和符号差异。}
